\documentclass[a4paper,11pt]{article}

\usepackage{kpfonts}
\usepackage[utf8]{inputenc}
\usepackage[margin=3cm]{geometry}

\usepackage{graphicx}
\usepackage{subcaption}
\usepackage{amsmath, amssymb, amsthm}
\usepackage{mathtools}
\DeclarePairedDelimiter\ceil{\lceil}{\rceil}
\DeclarePairedDelimiter\floor{\lfloor}{\rfloor}
\usepackage{microtype}
\usepackage{booktabs}
\usepackage{enumerate,paralist}
\usepackage{xspace}
\usepackage{comment}
\usepackage{xcolor}
\usepackage{array}
\usepackage[nocompress]{cite}
\usepackage[pdfpagelabels,colorlinks,allcolors=blue]{hyperref}
\usepackage[nameinlink,capitalize,noabbrev,sort]{cleveref}

\theoremstyle{plain}
\newtheorem{definition}{Definition}
\newtheorem{theorem}{Theorem}
\newtheorem{lemma}{Lemma}
\newtheorem{corollary}{Corollary}
\newtheorem{claim}{Claim}
\crefname{claim}{Claim}{Claims}
\Crefname{claim}{Claim}{Claims}
\newtheorem{proposition}{Proposition}
\newtheorem{observation}{Observation}
\theoremstyle{definition}  % switches off use of italics in following env's

\title{Smallest Cubic Non-1-Planar Graph}

%\author{Sergey Pupyrev}{Menlo Park, CA, USA}{spupyrev@gmail.com}{https://orcid.org/0000-0003-4089-673X}{}

%\authorrunning{S. Pupyrev}
%
%\Copyright{Sergey Pupyrev}
%
%\ccsdesc[500]{Theory of computation~Design and analysis of algorithms}
%\ccsdesc[500]{Mathematics of computing~Graph theory}
%\keywords{beyond planarity, 1-planar graph, cubic graph}

%\acknowledgements{}
%
%\hideLIPIcs
%\nolinenumbers

\graphicspath{{pics/}}

\definecolor{defblue}{rgb}{0.121,0.47,0.705}
\definecolor{defred}{rgb}{0.7,0.37,0.37}
\definecolor{defgreen}{rgb}{0.01,0.65,0.20}

\newcommand{\df}[1]{\textcolor{defblue}{\emph{#1}}}
\newcolumntype{C}[1]{>{\centering\arraybackslash}p{#1}}
\newcommand{\yes}{{\it \textcolor{defgreen}{yes}}}
\newcommand{\no}{{\it \textcolor{defred}{no}}}

\newcommand{\oops}{\texttt{OOPS}\xspace}

\begin{document}

\author{
    Sergey Pupyrev \\
    spupyrev@gmail.com
}

\date{}
\maketitle

\begin{abstract}
    A graph
    is 1-planar if it can be drawn in the plane such that each edge is crossed at most
    once---a natural generalization of planar graphs that has received increasing
    attention in graph drawing and beyond-planar graph theory.

    We search for the smallest cubic non-1-planar graph.
\end{abstract}

\section{Introduction}

Arguably among the most natural applications of a tool for recognizing 1-planar graphs is
finding the smallest non-1-planar instance in a family of graphs.
A particularly challenging variant of the problem is to
find the smallest 3-regular (cubic) graph that is not 1-planar. While it is easy to show
(see e.g., \cite{DKMW08}) that large random cubic graphs are not 1-planar, the
smallest such instance is unknown.

Throughout the paper, all graphs are finite, simple, and undirected.  The
\emph{order} of a graph $G$ is $|V(G)|$, and we write $n=|V(G)|$.  A graph is
\emph{cubic} if every vertex has degree three, and it is \emph{subcubic} if
every vertex has degree at most three.  A connected graph is
\emph{biconnected} if it has no cut vertex.  A graph is \emph{planar} if it
has a drawing without crossings.  The \emph{girth} of a graph containing a
cycle is the length of its shortest cycle.

\newtheorem*{problem1b}{Problem 1}
\begin{problem1b}[\cite{11011110_cubic,mathse_cubic,oops}]
    What is the smallest cubic non-1-planar graph?
\end{problem1b}

In an attempt to answer the question, Eppstein~\cite{11011110_cubic} manually constructed 1-planar
drawings of several graphs and suggested that either the Coxeter graph or the Tutte--Coxeter graph
is not 1-planar. As pointed out in \cite{oops}, the former is indeed 1-planar.
The latter is conjectured to be non-1-planar in \cite{oops}.

The paper consists of two parts. In \cref{sect:ub} we find two graphs of
order $n=30$ that are non-1-planar. In \cref{sect:lb} we verify that
all cubic graphs of order $n \le 28$ are 1-planar. Note that cubic graphs
necessarily have even number of vertices.
Most of the proofs in the paper are computational and verify 1-planarity
via a recently developed Optimized-One-Planarity-Solver (\oops)~\cite{oops} that
relies on encoding the problem using SAT.

\section{Two Cubic Graphs of Order 30 are Non-1-Planar}
\label{sect:ub}

This section will be filled out later.

\section{All Cubic Graphs of Order at most 28 are 1-Planar}
\label{sect:lb}

Our goal is to prove that every simple cubic graph with at most 28 vertices
is 1-planar.  Testing all graphs directly is impractical: there are
$2{,}094{,}480{,}864$ connected cubic graphs with $n=26$ and
$40{,}497{,}138{,}011$ with $n=28$.  Even at ten milliseconds per graph
(which one can achieve with \oops for \emph{simplest} instances),
the latter family alone would require more than twelve CPU-years.
\cref{table:cubic} and \cref{table:subcubic} list some families
of cubic graphs whose 1-planarity was verified in \cite{oops}, along with
the runtimes on a multi-core machine.
For example, the total computation of all cubic graphs with up to $n=24$ vertices and all cubic bipartite graphs with up to $n=28$ vertices combined
took $\approx$$1.5$ machine-months, averaging less than a second per instance.

\begin{table}[!t]
    \small
    \centering
    \caption{Computational results for cubic graphs.  The reported elapsed
    times were obtained by parallel computation.}
    \label{table:cubic}
    \begin{tabular}{|r|r|r|r||C{50pt}|}
        \toprule
        graph class & total count & total runtime & runtime per instance & all $1$-planar? \\
        \midrule
        $n=22$                    & $7,\!319,\!447$ & 5h 20m & 85 ms & \yes \\
        $n=24$                    & $117,\!940,\!535$ & 6d 9h 40m & $150$ ms & \yes \\
        \cmidrule{1-1}
        bipartite $n=26$             & $245,\!627$ & 25m & $200$ ms & \yes \\
        %girth$\ge 6$ $n=26$          & $181,\!227$ & $181,\!227$ & 4130 sec & $700$ ms & \yes \\
        girth$\ge 5$ $n=26$          & $31,\!478,\!584$ & 3d 6h 20m & $290$ ms & \yes \\
        $n=26$                    & $2,\!094,\!480,\!864$ & & & ? \\
        \cmidrule{1-1}
        bipartite $n=28$          & $2,\!291,\!589$ & 6h 55m & $340$ ms & \yes \\
        girth$\ge 6$ $n=28$          & $4,\!624,\!501$ & 2d 23h 20m & $1700$ ms & \yes \\
        $n=28$                    & $40,\!497,\!138,\!011$ & & & ? \\
        \cmidrule{1-1}
        bipartite $n=30$          & $23,\!466,\!857$ & & & ? \\
        \bottomrule
    \end{tabular}
\end{table}

TODO: merge two tables

\begin{table}[!t]
    \small
    \centering
    \caption{Computational results for selected subcubic graph families.}
    \label{table:subcubic}
    \begin{tabular}{@{}r p{205pt} r C{55pt}@{}}
        \toprule
        graph size & graph family & total count & all $1$-planar? \\
        \midrule
        $n=18$ & connected cubic                         & $41,\!301$ & \yes \\
        $n=18$ & biconnected, nonplanar, triangle-free cubic & $7,\!734$ & ? \\
        \cmidrule{1-1}
        $n=20$ & connected cubic                         & $510,\!489$ & \yes \\
        $n=20$ & biconnected, nonplanar, triangle-free cubic & $97,\!141$ & ? \\
        \cmidrule{1-1}
        $n=22$ & connected cubic                         & $7,\!319,\!447$ & \yes \\
        $n=22$ & connected, triangle-free cubic          & $1,\!435,\!720$ & ? \\
        $n=22$ & biconnected, nonplanar, triangle-free cubic & $1,\!432,\!712$ & ? \\
        \cmidrule{1-1}
        $n=24$ & connected cubic                         & $117,\!940,\!535$ & \yes \\
        $n=24$ & biconnected, nonplanar, triangle-free cubic & $23,\!751,\!969$ & ? \\
        \bottomrule
    \end{tabular}
\end{table}

\subsection{Building Blocks}

Our goal is to find the \emph{smallest subcubic non-1-planar} graphs.  We order
graphs lexicographically by $(|V|,|E|)$ and let $G$ be a minimum non-1-planar
graph of maximum degree at most three.  Thus every smaller subcubic graph is
1-planar. TODO: define \df{smallest}

Our first goal is to show that a smallest subcubic non-1-planar graph is biconnected, cubic, and triangle-free.

\begin{lemma}
    \label{lm:bi}
    A smallest subcubic non-1-planar graph is biconnected.
\end{lemma}

\begin{proof}
    Suppose, for a contradiction, that $G$ is a smallest subcubic
    non-1-planar graph and that $G$ is not biconnected.  If $G$ is
    disconnected, then each of its components has smaller order than $G$ and
    is therefore 1-planar.  Drawing these components in disjoint regions
    produces a 1-planar drawing of $G$, a contradiction.
    Hence $G$ is connected and has a cut vertex.  Every block of $G$ is a
    subcubic graph of smaller order and is therefore 1-planar.  The drawings
    of the blocks can be placed in disjoint wedges around their common cut
    vertices, producing a 1-planar drawing of $G$.  This is again a
    contradiction.
\end{proof}

\begin{lemma}
    \label{lm:cubic}
    A smallest subcubic non-1-planar graph is cubic.
\end{lemma}

\begin{proof}
    Suppose, for a contradiction, that $G$ is a smallest subcubic
    non-1-planar graph and that $G$ is not cubic. By \cref{lm:bi},
    $G$ is biconnected and therefore has minimum degree at least two.  Since
    $G$ is subcubic but not cubic, it has a vertex $v$ of degree two.  Let
    $u$ and $w$ be the neighbors of $v$.  Consider three cases.
    \begin{enumerate}[(a)]
        \item If $uw\notin E$, remove $v$ and insert $uw$.  The result $H$ is
        a smaller subcubic graph and hence is 1-planar.  In a drawing of $H$,
        replace the curve for $uw$ by the path $(u,v,w)$.  At most one path
        edge inherits the possible crossing of $uw$, so this produces a
        1-planar drawing of $G$, a contradiction.

        \item Suppose $uw\in E$ and $u,w$ have another common neighbor $x$.
        If $V(G)=\{u,v,w,x\}$, then $G$ is planar (the diamond graph, $K_4-e$), contradicting the choice
        of $G$.  Otherwise, the vertices $u,v,w$ have no neighbors outside
        $\{u,v,w,x\}$, so $x$ is a cut vertex of $G$.  This contradicts the
        preceding lemma.

        \item Suppose $uw\in E$ and $u,w$ have no other common neighbor.
        Contract the triangle $(u,v,w)$ to obtain a simple subcubic graph
        $H$ of smaller order.  In a 1-planar drawing of $H$, expand the
        contracted vertex inside a small disk into $(u,v,w)$.  This gives a
        1-planar drawing of $G$, a contradiction.
    \end{enumerate}
\end{proof}

\begin{comment}
A \df{planar two-terminal patch} is a planar graph $P$ with two designated
vertices $a,b$ that lie on a common face in some planar embedding of $P$.
To replace an edge $xy$ by $P$, delete $xy$, insert a disjoint copy of $P$,
and add the edges $xa$ and $yb$.

\begin{lemma}\label{lem:two-terminal-replacement}
    Let $G$ be obtained from $H$ by replacing an edge $xy$ with a planar
    two-terminal patch.  A 1-planar drawing of $H$ can be transformed
    into a 1-planar drawing of $G$ in which every edge of
    $E(H)\setminus\{xy\}$ retains its crossing status.  If $xy$ is uncrossed,
    then the entire replacement is uncrossed.
\end{lemma}

\begin{proof}
    Let $P$ be the inserted patch, with terminals $a$ and $b$.  Fix a
    1-planar drawing of $H$ and remove the curve representing $xy$.  Place a
    planar drawing of $P$ near $x$, using a common face of $a$ and $b$ as its
    outer face.  Join $x$ to $a$ locally and route $yb$ along the removed
    curve.  Thus $yb$ inherits the possible crossing of $xy$, while all other
    edges of the replacement are uncrossed.
\end{proof}
\end{comment}

\begin{lemma}
    \label{lem:minimal-triangle-free}
    A smallest subcubic non-1-planar graph is triangle-free.
\end{lemma}

\begin{proof}
    Suppose, for a contradiction, that $G$ is a smallest subcubic
    non-1-planar graph and that $G$ contains a triangle $T=(u,v,w)$.  By
    \cref{lm:bi,lm:cubic}, $G$ is biconnected and cubic.  Hence each vertex
    of $T$ has exactly one neighbor outside $T$; denote these neighbors by
    $u',v',w'$, and consider three cases.

    \begin{enumerate}[(a)]
        \item Suppose $u',v',w'$ are pairwise distinct.  Contract $T$ to a
        single vertex.  The resulting graph is simple, subcubic, and of
        smaller order, so it is 1-planar.  Expanding the contracted vertex
        inside a small disk gives a 1-planar drawing of $G$, a contradiction.

        \item Suppose $u'=v'=w'$.  Then connectedness and cubicity imply that
        $G=K_4$, contradicting the non-1-planarity of $G$.

        \item Otherwise, after relabeling, $u'=v'=x$ and $w'=y\ne x$.  The
        vertices $u,v,w,x$ induce a diamond.  Let $p$ be the third neighbor
        of $x$.  If $p=y$, then the diamond is attached to the rest of $G$
        only through $y$, so $y$ is a cut vertex, contradicting
        biconnectivity.

        Hence $p\ne y$.  Contract the diamond to a single vertex $z$.  The
        resulting graph $H$ is simple, subcubic, and of smaller order, so it
        is 1-planar.  In a 1-planar drawing of $H$, expand $z$ inside a small
        disk into the diamond, attaching $zp$ at $x$ and $zy$ at $w$.  This
        gives a 1-planar drawing of $G$, a contradiction.
    \end{enumerate}
\end{proof}

\Cref{lm:bi,lm:cubic,lem:minimal-triangle-free} imply the following.

\begin{corollary}\label{cor:minimal-subcubic}
    A smallest subcubic non-1-planar graph is biconnected, cubic, and
    triangle-free.
\end{corollary}

We next strengthen \cref{cor:minimal-subcubic} by prescribing edges that must
remain uncrossed.  A graph $G$ is \df{$k$-flexible} if, for every set
$F\subseteq E(G)$ with $|F|\leq k$, it has a 1-planar drawing in which every
edge of $F$ is uncrossed.  In particular, $0$-flexibility is equivalent to
1-planarity.

\begin{lemma}
    \label{lem:flexibility-core}
    For every integer $k\geq 0$, a smallest subcubic graph that is not
    $k$-flexible is biconnected, cubic, and triangle-free.
\end{lemma}

\begin{proof}
    Let $G$ be such a graph.  For any set $F\subseteq E(G)$ with $|F|\leq k$,
    distribute the edges of $F$ among the components or blocks of $G$.  Every
    proper component or block is smaller than $G$ and is therefore
    $k$-flexible.  Its drawing may consequently be chosen with its part of
    $F$ uncrossed.  Placing component drawings in disjoint regions, or block
    drawings in disjoint wedges at their common cut vertices, gives the
    required drawing of $G$.  Hence $G$ is biconnected.

    A biconnected graph with a vertex of degree at most one is either a single
    vertex or a single edge, and is therefore $k$-flexible.  It remains to
    exclude a vertex $v$ of degree two.  Let its neighbors be
    $u$ and $w$.  If $uw\notin E(G)$, delete $v$ and add $uw$.  Retain every
    prescribed edge outside $uv,vw$, and prescribe $uw$ if $uv$ or $vw$
    belongs to $F$.  Subdividing an uncrossed drawing of $uw$ restores both
    $uv$ and $vw$ without crossings.

    Suppose instead that $uw\in E(G)$.  If $u$ and $w$ have another common
    neighbor $x$, then either $G$ is planar or $x$ is a cut vertex, contrary
    to the preceding paragraph.  Otherwise identify $u,v,w$ to a single
    vertex.  The resulting graph is simple and subcubic.  Retain prescribed
    edges outside the triangle, map its prescribed attachment edges to their
    images, and omit its prescribed internal edges.  Expanding the identified
    vertex inside a disk restores the triangle, keeps its internal edges
    uncrossed, and preserves the crossing status of every mapped attachment.
    Thus $G$ is cubic.

    Finally, let $T$ be a triangle.  If its three external neighbors are
    distinct, contract $T$ to $z$.  Map a prescribed attachment $v_ix_i$ to
    $zx_i$, retain prescribed outside edges, and omit prescribed triangle
    edges from the mapped set; the map does not increase its size.  Expanding
    an uncrossed mapped attachment and drawing the triangle inside the vertex
    disk preserves every prescribed edge.  In the repeated-neighbor cases,
    apply the iterative two-terminal replacement used in the proof of
    \cref{lem:minimal-triangle-free}.  The prescribed-edge construction in
    \cref{lem:two-terminal-replacement} does not increase the number of
    prescribed edges.  Every case therefore reduces to a smaller subcubic
    graph, a contradiction.
\end{proof}

\begin{lemma}[4-cycle expansion]\label{lem:4-cycle-expansion}
    Let $G$ be a triangle-free cubic graph containing the cycle
    $C=(v_0,v_1,v_2,v_3)$, and let $x_i$ be the neighbor of $v_i$ outside
    $C$.  Form $H$ by deleting $C$, adding adjacent vertices $a,b$, joining
    $a$ to $x_0,x_1$, and joining $b$ to $x_2,x_3$.  Write
    $c_i=v_iv_{i+1}$ and $s_i=v_ix_i$, with indices modulo four.
    \begin{enumerate}[(i)]
        \item If $H$ has a 1-planar drawing in which $ab$ is uncrossed, then
        $G$ is 1-planar.
        \item If $H$ has a 1-planar drawing in which $ab$ and the image of
        $s_1$ are uncrossed, then $G$ has a 1-planar drawing in which
        $c_0,c_1,c_2$ are uncrossed.
    \end{enumerate}
\end{lemma}

\begin{proof}
    Triangle-freeness implies $x_0\ne x_1$ and $x_2\ne x_3$, so $H$ is a
    simple cubic graph; equality between opposite external neighbors does not
    create parallel edges.  Take a narrow disk containing the uncrossed curve
    $ab$ and small neighborhoods of its endpoints.  Inside the disk, split
    $a$ into $v_0,v_1$, split $b$ into $v_2,v_3$, and preserve the four
    external incidences.  Up to reflection, the four possible boundary orders
    are
    \[
       0123,\quad 0132,\quad 0231,\quad 0321.
    \]
    For part (i), the first and last orders require no crossing, while the
    other two use the single crossing $c_1\mathbin\times c_3$.  For part (ii),
    the first and last again require no crossing, while the other two use
    $c_3\mathbin\times s_1$.  These cases exhaust the endpoint rotations.  In
    every case the crossing pairs form a matching, their planarizations were
    verified to be planar with the prescribed boundary order, and the edges
    required in the corresponding part remain uncrossed.
\end{proof}

The analogous construction pairing $x_1,x_2$ at $a$ and $x_3,x_0$ at $b$
gives the corresponding second expansion.

\begin{lemma}[4-cycle reduction through order 24]\label{lem:order24-reduction}
    Suppose every subcubic graph of order at most 22 is $3$-flexible and
    every biconnected, nonplanar cubic graph of order 24 and girth at least
    five is $2$-flexible.  Then every cubic graph of order at most 24 is
    $2$-flexible.
\end{lemma}

\begin{proof}
    By \cref{lem:flexibility-core}, it is enough to consider a biconnected
    triangle-free cubic graph $G$ of order 24 and a prescribed set $F$ of at
    most two edges, since a smaller counterexample is covered by the first
    hypothesis.  We may assume that $G$ is nonplanar.  If $G$ has no 4-cycle,
    the second hypothesis applies.  Otherwise choose a 4-cycle $C$.

    First suppose that $F$ is not a pair of adjacent edges of $C$.  Choose
    one of the two 4-cycle reductions so that every prescribed cycle
    edge is created inside an endpoint disk: one split handles the opposite
    pair $c_0,c_2$, and the other handles $c_1,c_3$.  In the resulting
    order-22 graph prescribe the new middle edge and the images of all edges
    of $F$ outside $C$.  This is a set of at most three edges.  A drawing
    supplied by $3$-flexibility expands by \cref{lem:4-cycle-expansion};
    prescribed outside edges are unchanged, and prescribed cycle edges are
    drawn uncrossed inside the endpoint disks.

    The remaining case is $F=\{c_i,c_{i+1}\}$.  Rotate the labels so that
    these are $c_0,c_1$, use the displayed split, and prescribe only $ab$ and
    the image of $s_1$.  By part (ii) of \cref{lem:4-cycle-expansion}, the
    resulting $3$-flexible order-22 drawing expands while keeping both
    prescribed edges uncrossed.
\end{proof}

Consequently, it suffices to verify $3$-flexibility for biconnected,
triangle-free cubic graphs of order at most 22 and $2$-flexibility for
biconnected, nonplanar cubic graphs of order 24 and girth at least five.

\begin{lemma}[5-cycle-to-path expansion]\label{lem:5-cycle-path-expansion}
    Let $C=(v_0,\ldots,v_4)$ be a chordless 5-cycle in a cubic graph $G$,
    and let $x_i$ be the neighbor of $v_i$ outside $C$.  Suppose that the
    vertices $x_0,\ldots,x_4$ are distinct.  Write $c_i=v_iv_{i+1}$ and
    $s_i=v_ix_i$, with indices modulo five.  Form $H$ by deleting $C$,
    introducing a path $abc$, and adding the edges
    $ax_0,ax_1,bx_2,cx_3,cx_4$.  If $H$ has a 1-planar drawing in which the
    three edges incident with $b$ are uncrossed, then $G$ is 1-planar.
\end{lemma}

\begin{proof}
    Take a regular neighborhood of the uncrossed path $abc$.  Up to
    reflection, the five attachment portions have the four boundary orders
    shown in \cref{tab:5-cycle-path-expansion}.  For each order, the table
    lists a matching of crossing pairs whose planarization was verified to be
    planar with the prescribed boundary order.  The only old attachment used
    in a local crossing is $s_2=v_2x_2$, which is uncrossed outside the
    neighborhood.  Replacing the path by the corresponding local drawing of
    the 5-cycle therefore gives a 1-planar drawing of $G$.
\end{proof}

\begin{table}[!ht]
    \centering
    \caption{Crossing pairs for the 5-cycle-to-path expansion.}
    \label{tab:5-cycle-path-expansion}
    \begin{tabular}{@{}>{$}c<{$}@{\qquad}>{$}l<{$}@{}}
        \toprule
        \multicolumn{1}{c}{boundary order} &
        \multicolumn{1}{c}{crossing pairs} \\
        \midrule
        01234&\varnothing\\
        01243&c_2\mathbin\times c_4\\
        01342&c_4\mathbin\times s_2\\
        01432&c_1\mathbin\times c_4\\
        \bottomrule
    \end{tabular}
\end{table}

\begin{lemma}[Overlapping 5-cycle expansion]
    \label{lem:overlapping-5-cycle-expansion}
    Let $C=(v_0,\ldots,v_4)$ be a 5-cycle in a cubic graph $G$ of girth five,
    and let $x_i$ be the neighbor of $v_i$ outside $C$.  Suppose that the
    vertices $x_0,\ldots,x_4$ are distinct.  Form $H$ by deleting $C$,
    introducing a path $abc$, and adding the edges
    $ax_0,ax_1,bx_2,cx_3,cx_4$.  Suppose that $abc$ lies on another 5-cycle
    $D$ of $H$.  Replace $D$ by a three-vertex path, assigning the attachment
    at $b$ to its middle vertex, and call the resulting order-24 graph $K$.
    Suppose that this reduction is simple and connected.  If $K$ has a
    1-planar drawing in which the two edges of the new path are uncrossed,
    then $G$ is 1-planar.
\end{lemma}

\begin{proof}
    The two vertices of $D-\{a,b,c\}$ are outside the original 5-cycle.
    Girth five excludes the case in which they attach at consecutive
    vertices of the original cycle.  Up to a dihedral relabeling, the patch
    to be restored therefore consists of
    $v_0,\ldots,v_4,x_0,x_3$, with the additional edges
    $v_0x_0,v_3x_3,x_0x_3$.

    Take a regular neighborhood of the uncrossed path in $K$.  Its five
    attachment portions occur, up to reflection, in the four boundary orders
    listed in \cref{tab:overlapping-5-cycle-expansion}.  The local drawings in
    \cref{tab:overlapping-5-cycle-expansion} restore the seven-vertex patch.
    Here $d_0=v_0x_0$ and $d_3=v_3x_3$.  Every displayed crossing uses two
    edges of the restored patch, so no further edge of the drawing of $K$
    must be kept uncrossed.  The four rows exhaust the boundary orders.  In
    every row the crossing pairs form a matching, and their planarizations
    were verified to be planar with the prescribed boundary order.
\end{proof}

\begin{table}[!ht]
    \centering
    \caption{Crossing pairs for the overlapping 5-cycle expansion.}
    \label{tab:overlapping-5-cycle-expansion}
    \begin{tabular}{@{}>{$}c<{$}@{\qquad}>{$}l<{$}@{}}
        \toprule
        \multicolumn{1}{c}{boundary order} &
        \multicolumn{1}{c}{crossing pairs} \\
        \midrule
        01234&c_4\mathbin\times d_3\\
        01243&\{c_2\mathbin\times c_4,d_0\mathbin\times d_3\}\\
        01342&c_3\mathbin\times d_0\\
        01432&c_2\mathbin\times d_0\\
        \bottomrule
    \end{tabular}
\end{table}

\begin{lemma}[Contracted 5-cycle expansion]\label{lem:5-cycle-star-expansion}
    Let $C=(v_0,\ldots,v_4)$ be a chordless 5-cycle in a cubic graph $G$,
    and let $x_i$ be the neighbor of $v_i$ outside $C$.  Suppose that the
    vertices $x_0,\ldots,x_4$ are distinct, and let $J$ be obtained by
    contracting $C$ to one vertex.  Write $c_i=v_iv_{i+1}$ and
    $s_i=v_ix_i$, with indices modulo five, and call the edges $s_i$ the
    attachments of $C$.
    \begin{enumerate}[(i)]
        \item If $J$ has a 1-planar drawing in which all five edges
        $s_0,\ldots,s_4$ are uncrossed, then, for every $i$, the graph $G$
        has a 1-planar drawing in which the three edges incident with $v_i$
        are uncrossed.
        \item Let $e\in E(J)$.  If $e$ is not an attachment, suppose that
        $J$ has a 1-planar drawing in which $e$ and at least three
        attachments are uncrossed.  If $e=s_i$, instead suppose that $e$ and
        at least three of the other four attachments are uncrossed.  In
        either case, $G$ has a 1-planar drawing in which $e$ is uncrossed.
    \end{enumerate}
\end{lemma}

\begin{proof}
    Take a disk around the contracted vertex.  Fixing $s_0$ as the first
    attachment and identifying an order with its reversal leaves twelve
    possible cyclic orders on the boundary of the disk.  These are exactly
    the rows of \cref{tab:5-cycle-star-expansion}.  In every row, the crossing
    pairs in each column form a matching, avoid the edges required to remain
    uncrossed, and yield a planar graph after planarization with the prescribed
    boundary order.  The second column
    of \cref{tab:5-cycle-star-expansion} gives a local drawing for each order
    whose attachment crossings use only $s_0,s_1$.  Relabeling the cycle maps
    any chosen consecutive pair $s_j,s_{j+1}$ to this pair.  Any three
    attachments contain a consecutive pair.  If $e$ is an attachment, any
    three of the other four contain a consecutive pair avoiding $e$.
    Choosing this pair ensures that the local drawing does not cross $e$;
    if $e$ is not an attachment, it lies outside the replacement disk.
    This proves (ii).

    For (i), it is enough by symmetry to preserve the star
    $\{c_4,c_0,s_0\}$.  The third column gives a local drawing for every
    boundary order in which these three edges are uncrossed.  Its attachment
    crossings use only $s_1,s_2,s_3,s_4$, all of which are uncrossed outside
    the disk.  Thus every listed local drawing is 1-planar, and inserting it
    into the disk proves both statements.
\end{proof}

\begin{table}[!ht]
    \centering
    \caption{Crossing pairs for the 5-cycle star expansion.  The ordinary
    column permits local crossings on $s_0,s_1$; the last column preserves
    the star $\{c_4,c_0,s_0\}$.}
    \label{tab:5-cycle-star-expansion}
    \begin{tabular}{@{}>{$}c<{$}@{\qquad}>{$}l<{$}@{\qquad}>{$}l<{$}@{}}
        \toprule
        \multicolumn{1}{c}{boundary order} &
        \multicolumn{1}{c}{ordinary expansion} &
        \multicolumn{1}{c}{star at $v_0$} \\
        \midrule
        01234&\varnothing&\varnothing\\
        01243&c_2\mathbin\times c_4&c_2\mathbin\times s_4\\
        01324&c_1\mathbin\times c_3&c_1\mathbin\times c_3\\
        01342&\{c_1\mathbin\times s_0,c_2\mathbin\times c_4\}&
               \{c_1\mathbin\times c_3,s_2\mathbin\times s_4\}\\
        01423&\{c_1\mathbin\times c_3,c_4\mathbin\times s_1\}&
               c_1\mathbin\times s_4\\
        01432&c_1\mathbin\times c_4&
               \{c_1\mathbin\times s_4,c_3\mathbin\times s_2\}\\
        02134&c_0\mathbin\times c_2&c_2\mathbin\times s_1\\
        02143&c_2\mathbin\times s_0&
               \{c_2\mathbin\times s_1,s_3\mathbin\times s_4\}\\
        02314&c_3\mathbin\times s_1&c_3\mathbin\times s_1\\
        02413&\{c_1\mathbin\times c_4,c_2\mathbin\times s_0,
                 c_3\mathbin\times s_1\}&
               \{c_1\mathbin\times c_3,s_1\mathbin\times s_3,
                 s_2\mathbin\times s_4\}\\
        03124&\{c_0\mathbin\times c_2,c_3\mathbin\times s_0\}&
               \{c_1\mathbin\times c_3,s_1\mathbin\times s_3\}\\
        03214&c_0\mathbin\times c_3&
               \{c_1\mathbin\times s_3,c_3\mathbin\times s_1\}\\
        \bottomrule
    \end{tabular}
\end{table}

\subsection{Verification}

The reductions above leave five computational claims.  All enumerations below
are exhaustive and contain one representative of each isomorphism class.

We first describe the mathematical objects checked in these computations.  A
\emph{crossing pattern} in a graph $G$ is a set $M$ of unordered pairs of
nonadjacent edges such that no edge occurs in more than one pair.  Let
$\operatorname{supp}(M)$ be the set of edges occurring in the pairs of $M$,
and let $G_M$ be obtained by replacing each pair $\{uv,xy\}\in M$ with a new
vertex adjacent to $u,v,x,y$.

If $G_M$ is planar, then $G$ is 1-planar.  Indeed, consider a planar embedding
of $G_M$ and a small disk around each new vertex.  If the two portions of each
original edge alternate around the new vertex, replace that vertex by a
crossing; otherwise join the two portions of each original edge without a
crossing.  Since no edge belongs to two pairs of $M$, every edge is crossed at
most once.  Moreover, every edge outside $\operatorname{supp}(M)$ is
uncrossed.  The computational proofs below consist of finding such crossing
patterns and testing the planarity of the corresponding graphs $G_M$.  For a
planar graph, the empty crossing pattern suffices.

\begin{claim}\label{clm:flexible22}
    Every biconnected, nonplanar, triangle-free cubic graph with
    $n\leq 22$ is $3$-flexible.
\end{claim}

\begin{proof}[Proof (computational)]
    This family contains $1{,}538{,}493$ isomorphism classes.  For every
    graph $G$ in the enumeration and every set
    $F\subseteq E(G)$ with $|F|\leq 3$, the computation found a crossing
    pattern $M$ such that $G_M$ is planar and
    $F\cap\operatorname{supp}(M)=\varnothing$.  Hence $G$ has a 1-planar
    drawing in which every edge of $F$ is uncrossed.  The complete computation
    took 17 hours 49 minutes on the 48-core machine.
\end{proof}

\begin{claim}\label{clm:flexible24}
    Every biconnected, nonplanar cubic graph with $n=24$ and girth $\geq 5$
    is $2$-flexible.
\end{claim}

\begin{proof}[Proof (computational)]
    This family contains $1{,}620{,}470$ isomorphism classes.  For every graph
    $G$ in the enumeration and every set $F\subseteq E(G)$ with $|F|\leq 2$,
    the computation found a crossing pattern $M$ such that $G_M$ is planar and
    $F\cap\operatorname{supp}(M)=\varnothing$.  Thus every edge of $F$ is
    uncrossed in the resulting 1-planar drawing.
\end{proof}

\begin{claim}\label{clm:flexible26}
    Every biconnected, nonplanar cubic graph with $n=26$ and girth
    at least five is $1$-flexible.
\end{claim}

\begin{proof}[Proof (computational)]
    The filtered input contains $31{,}478{,}465$ isomorphism classes.  Of
    these, $31{,}297{,}238$ have girth five and $181{,}227$ have larger girth.
    First consider the graphs of girth five.  For each such graph $G$, the
    computation contracted every 5-cycle, canonically labeled the resulting
    order-22 graph, and retained the lexicographically least code.  Equal
    retained graphs were verified only once.  Write $J$ for one retained
    graph and $s_0,\ldots,s_4$ for the five edges incident with its
    degree-five vertex.  This produced $4{,}700{,}649$ distinct cores.

    For every core completed by the first pass, the computation found a
    crossing pattern with all five $s_i$ uncrossed.  In addition, for every
    edge $e\in E(J)$ it found a crossing pattern in which $e$ and at least
    three suitable attachments are uncrossed: all five attachments are
    eligible when $e$ is not one of them, and the other four are eligible
    when $e$ is an attachment.  Each resulting planarization was checked.
    These conditions do not depend on the cyclic order at the contracted
    vertex.  An edge on the contracted 5-cycle is preserved using part (i)
    of \cref{lem:5-cycle-star-expansion}; every other edge is preserved using
    part (ii).

    A 120-second limit retains the cores not completed by this first pass.
    For each retained core, the second pass expands its degree-five vertex in
    all \(4!/2=12\) cyclic orders.  It keeps the biconnected, nonplanar
    order-26 expansions of girth at least five and verifies their
    \(1\)-flexibility directly.  Every original parent represented by the
    retained core is one of these expansions.  Hence the two passes together
    cover every girth-five graph.

    The remaining order-26 graphs have girth at least six.  For every such
    graph $G$ and every edge $e\in E(G)$, the computation directly found a
    crossing pattern whose planarization is planar and whose support avoids
    $e$.  Any such graph retained by the first pass is retried
    directly in the second pass.

    Together with the $181{,}227$ unchanged graphs of larger girth, the first
    pass processed $4{,}881{,}876$ records.  It completed $4{,}881{,}171$ of
    them and retained 705 cores.  The second pass expanded these cores into
    $8{,}460$ order-26 parents and verified every parent as $1$-flexible.  No
    record remained unresolved.  The two solver passes took 22.97 hours and
    26 minutes, respectively, on the 48-core machine.
\end{proof}

\begin{claim}\label{clm:girth5-28}
    Every biconnected, nonplanar cubic graph with $n=28$ and girth $=5$ is
    1-planar.
\end{claim}

\begin{proof}[Proof (computational)]
    The exhaustive enumeration contains $652{,}157{,}758$ isomorphism
    classes.  For every graph $G$ in this family, the computation enumerates
    the 5-cycle-to-path reductions $(H,b)$ defined in
    \cref{lem:5-cycle-path-expansion}.  Girth five makes every first
    reduction simple.  Write the replacement path as $a b c$.

    If the replacement path lies on a second 5-cycle, the computation
    performs the second path reduction with $b$ assigned to its middle
    attachment.  When the resulting order-24 graph $K$ is biconnected and
    has girth at least five, it is either planar or covered by
    \cref{clm:flexible24}.  Thus its two path edges can be kept uncrossed, and
    \cref{lem:overlapping-5-cycle-expansion} proves that $G$ is 1-planar.

    Otherwise the computation finds one of two overlapping cycles.  With
    $L,M$ denoting attachments at $a,b$, respectively, their cyclic words are
    \[
        b\,a\,L\,O\,M
        \qquad\hbox{or}\qquad
        b\,a\,L\,O\,O\,M .
    \]
    Contract the complete union of the original 5-cycle and this overlap to
    one vertex.  The resulting quotient has order 21 and one vertex of
    degree six in the first case, or order 20 and one vertex of degree seven
    in the second case; every other vertex has degree three.  Equal
    canonical quotients are checked only once.

    A finite exhaustive check gives the following exact expansion rule.
    Every one of the 60 boundary orders
    of the degree-six patch expands.  For the degree-seven patch, 355 of the
    360 boundary orders expand; the five exceptions are
    \[
      0245136,\quad0246135,\quad0264153,\quad
      0315426,\quad0316425.
    \]
    The check surrounds the patch by an uncrossed wheel in the prescribed
    order and verifies the resulting crossing pattern and planarization.

    A degree-six quotient is certified by a drawing in which all six hub
    edges are uncrossed.  A degree-seven quotient is certified if, for every
    bijection between the seven patch attachments and the seven hub edges,
    there is such a drawing whose rotation avoids the five exceptional
    orders.  The rotation is obtained from the checked planarization and is
    therefore part of the certificate.  Replacing the hub by the appropriate
    local drawing then restores $G$.

    If no direct or hub certificate is found, the computation retains the
    lexicographically least canonically labeled reduction $(H,b)$ and seeks a
    drawing of $H$ in which the three edges incident with $b$ are uncrossed.
    Such a drawing restores $G$ by
    \cref{lem:5-cycle-path-expansion}.  Isomorphic hub and marked reductions
    are checked only once.

    Some of these constrained checks are substantially harder than ordinary
    1-planarity.  A fixed time limit was therefore used only to divide the
    exhaustive computation: for each unfinished reduction, every original
    graph assigned to that reduction was tested directly for 1-planarity.
    No conclusion was drawn from reaching the time limit.

    The classification assigned one record to each of the
    $652{,}157{,}758$ input graphs.  Of these, $254{,}354{,}420$ were covered
    by the first expansion argument using \cref{clm:flexible24}; the remaining
    $397{,}803{,}338$ were assigned a degree-six, degree-seven, or
    marked-reduction record.  After
    canonical labeling, the latter records gave $3{,}788{,}793$ distinct
    constrained instances.  The computation settled $3{,}784{,}587$ of
    them and retained $4{,}206$ at the time limit.  These retained records
    represented $1{,}508{,}639$ original graphs, each of which was verified
    directly as 1-planar.  Thus every graph in the enumeration is covered
    either by one of the expansion arguments above or by a directly checked
    drawing.
\end{proof}

\begin{claim}\label{clm:girth6-28}
    Every biconnected, nonplanar cubic graph with $n=28$ and girth $\geq 6$ is
    1-planar.
\end{claim}

\begin{proof}[Proof (computational)]
    This family contains $4{,}624{,}501$ isomorphism classes.  For every graph
    $G$ in the enumeration, the computation found a crossing pattern $M$ such
    that $G_M$ is planar.  Hence every graph in the family is 1-planar.
\end{proof}

The running times for the five computations are summarized in
\cref{tab:cubic28-runtime}.  All five computations are complete.

For Claim~3, preparation reduced the $31{,}478{,}465$ input graphs to
$4{,}881{,}876$ records.  The first solver pass took 22.97 hours and retained
705 hard cores.  Expanding these cores into $8{,}460$ concrete parents and
verifying them directly took another 26 minutes.  Thus Claim~3 took about
23 hours 24 minutes of measured solver time.

For Claim~4, classifying the $652{,}157{,}758$ input graphs took about 12 hours
on the 48-core machine.  The computation found $254{,}354{,}420$ direct
certificates, $350{,}383{,}009$ degree-six quotients, $45{,}970{,}275$
degree-seven quotients, and $1{,}450{,}054$ marked reductions.  Canonical
labeling reduced the latter three classes to $3{,}788{,}793$ distinct records.
With a 60-second limit, the certificate computation found 17 planar records,
verified $3{,}784{,}570$ further records, and retained $4{,}206$.  The retained
records represented $1{,}508{,}639$ original graphs, all of which were verified
directly as 1-planar.  The complete computation took 52.5 hours on the
48-core machine.

\begin{table}[!ht]
    \small
    \centering
    \caption{Measured wall-clock times on 48 cores.}
    \label{tab:cubic28-runtime}
    \begin{tabular}{@{}lrrc@{}}
        \toprule
        computation & instances & running time & verified \\
        \midrule
        \cref{clm:flexible22} & $1{,}538{,}493$ & 17 hours 49 minutes & yes \\
        \cref{clm:flexible24} & $1{,}620{,}470$ & 41 hours 21 minutes & yes \\
        \cref{clm:flexible26} & $4{,}881{,}876+8{,}460^\ast$ & 23 hours 24 minutes & yes \\
        \cref{clm:girth5-28}  & $3{,}788{,}793+1{,}508{,}639^\ast$ & 52.5 hours & yes \\
        \cref{clm:girth6-28}  & $4{,}624{,}501$ & 3 days & yes \\
        \bottomrule
    \end{tabular}
    \par\smallskip
    {\footnotesize $^\ast$First-pass records plus directly verified fallback parents.\par}
\end{table}

\begin{theorem}\label{thm:cubic28}
    Every simple cubic graph with at most 28 vertices is 1-planar.
\end{theorem}

\begin{proof}
    A planar graph is $k$-flexible for every $k$, so every obstruction to a
    flexibility statement below is nonplanar.  We first apply the flexibility
    results at the smaller orders.  If a subcubic graph of order at most 22
    were not $3$-flexible, a smallest such graph would be biconnected, cubic,
    and triangle-free by
    \cref{lem:flexibility-core}, contrary to \cref{clm:flexible22}.  Thus every
    subcubic graph of order at most 22 is $3$-flexible.  It follows from
    \cref{lem:order24-reduction,clm:flexible24} that every cubic graph of order
    at most 24 is $2$-flexible.  A smallest subcubic graph of order at most 24
    that is not $2$-flexible would again be cubic by
    \cref{lem:flexibility-core}; consequently, every subcubic graph of order
    at most 24 is $2$-flexible.

    We next show that every subcubic graph of order at most 26 is
    $1$-flexible.  Suppose otherwise, and choose a smallest counterexample
    $G$.  By \cref{lem:flexibility-core}, the graph $G$ is biconnected, cubic,
    and triangle-free.  The preceding paragraph excludes orders at most 24,
    and a cubic graph has even order; hence $G$ has order 26.  Fix an edge
    $e\in E(G)$.  If $G$ contains a 4-cycle, choose one of the
    two reductions in \cref{lem:4-cycle-expansion} so that, when $e$ belongs
    to the cycle, it is restored inside an endpoint disk.  In the resulting
    graph of order 24, prescribe the edge $ab$ and, when $e$ lies outside the
    cycle, its image.  Its $2$-flexibility gives a drawing that expands to a
    drawing of $G$ in which $e$ is uncrossed.  If $G$ has no 4-cycle, then it
    has girth at least five, so \cref{clm:flexible26} directly gives a
    1-planar drawing in which $e$ is uncrossed.  Since $e$ was arbitrary,
    both cases contradict the choice of $G$.

    Finally, suppose that a subcubic graph of order at most 28 is not
    1-planar, and choose a smallest such graph $G$.  By
    \cref{lem:flexibility-core} with $k=0$, the graph $G$ is biconnected,
    cubic, and triangle-free.  The preceding paragraph excludes orders at
    most 26, and cubic graphs have even order; hence $G$ has order 28.  If $G$
    contains a 4-cycle, apply
    \cref{lem:4-cycle-expansion}.  The resulting graph has order 26, and its
    $1$-flexibility provides a 1-planar drawing in which $ab$ is uncrossed;
    this drawing expands to a 1-planar drawing of $G$.  Otherwise $G$ has
    girth at least five.  If its girth is five, \cref{clm:girth5-28} applies;
    if its girth is at least six, \cref{clm:girth6-28} applies.  Every case is
    a contradiction.
\end{proof}





\bibliographystyle{plainurl}
\bibliography{cubic}

\end{document}
